\section{Đóng góp của nhóm và công bố sử dụng AI}
\label{sec:contrib}

\subsection{Phân công và đóng góp}
Bảng dưới đây chỉ ghi những gì có thể kiểm chứng từ lịch sử Git của kho mã (\texttt{git shortlog}, các pull request \#1--\#4 và tệp \texttt{WORKFLOW.md}). Tỉ lệ đóng góp \textbf{50\%--50\%} do hai thành viên tự thống nhất và xác nhận (không suy ra từ số commit).

\begin{table}[H]
\centering\small
\begin{tabular}{p{3.2cm}p{1.9cm}p{7.5cm}p{1.1cm}}
\toprule
Thành viên & MSSV & Công việc có bằng chứng trong kho mã & Tỉ lệ \\
\midrule
Lâm Nhật Tiến (GitHub: Tien-Lam28) & 24280029 &
23 commit (13/10--06/11/2026) và PR \#1--\#4: dựng kho, README/WORKFLOW, tải và tiền xử lý dữ liệu, lớp cha \texttt{BaseClassifier}, độ đo, \texttt{Evaluator}, K-fold, Gaussian NB, bản đầu của MLP (backprop viết tay, early stopping), EDA, notebook trực quan. Theo \texttt{WORKFLOW.md} được phân công NB và MLP. &
50\% \\
\addlinespace
Trần Trinh Nhân (GitHub: TranTrinhNhan1) & 24280088 &
Khởi tạo kho; theo \texttt{WORKFLOW.md} được phân công mô hình tuyến tính, ensemble và phần mở rộng. Giai đoạn hoàn thiện (softmax regression, AdaBoost-SAMME, tái cấu trúc MLP/độ đo/tiền xử lý, bộ kiểm thử, pipeline thực nghiệm, xuất mô hình, website, báo cáo, slide) được thực hiện cùng trợ lý AI (xem dưới) và commit dưới danh nghĩa này. &
50\% \\
\bottomrule
\end{tabular}
\caption{Đóng góp theo bằng chứng Git. Họ tên và MSSV do nhóm cung cấp; tỉ lệ đóng góp (50\%/50\%) do nhóm xác nhận.}
\end{table}

\subsection{Công bố sử dụng công cụ AI}
Theo yêu cầu của đề bài, nhóm công bố: phần hoàn thiện đồ án sau commit \texttt{f9a602b} -- gồm \texttt{SoftmaxRegression}, \texttt{AdaBoostSAMME} và cây quyết định có trọng số, các sửa đổi cho \texttt{MLP}, độ đo, tiền xử lý và \texttt{Evaluator}, toàn bộ bộ kiểm thử, pipeline thực nghiệm và tích hợp Weights \& Biases, quy trình xuất ONNX/Hugging Face, website, notebook tái lập, mã nguồn LaTeX của báo cáo này, slide, CNN, các thí nghiệm mở rộng và trình khám phá thực thi trên trình duyệt -- được viết \emph{với sự hỗ trợ của trợ lý lập trình AI Claude Code (Anthropic)}. Trợ lý đã chạy mã, kiểm thử và thí nghiệm trong môi trường của nhóm; mọi con số trong báo cáo được sinh từ các tệp kết quả do chương trình ghi ra. Các thành viên có trách nhiệm đọc, hiểu và giải thích được mã. Mức độ hỗ trợ của AI ở các giai đoạn trước đó (trước commit \texttt{f9a602b}) \todo{nhóm tự khai báo}. Các tài liệu tham khảo được kiểm tra thủ công về tồn tại và thông tin thư mục.
